\section{Архитектура программного комплекса «Модейлер»}

\subsection{Общее описание платформы}

«Модейлер» — кроссплатформенная программная система для интерактивного моделирования и симуляции роботизированных устройств. Платформа предназначена для инженеров-разработчиков, исследователей и специалистов по автоматизации, решающих задачи верификации кинематических и динамических характеристик роботов без физического прототипирования.

В области промышленной робототехники задачи, для которых применяется «Модейлер», охватывают проектирование роботизированных ячеек, отладку управляющих алгоритмов в режиме программного замкнутого контура (SIL), а также подготовку траекторий движения. В сфере исследовательских и образовательных применений платформа обеспечивает возможность экспериментов с моделями на основе описания роботов в формате URDF\footnote{Unified Robot Description Format — стандартный XML-формат, используемый в экосистеме ROS.} без необходимости развёртывания реального оборудования.

В отличие от универсальных систем физического моделирования (рассмотренных в разделе~1.2), «Модейлер» ориентирован именно на робототехнические применения: нативная поддержка URDF-моделей, интеграция с ROS~2 и специализированный GPU-ускоренный физический движок составляют его ключевые характеристики.

\subsection{Многоуровневая архитектура}

Архитектура «Модейлера» организована по принципу строгого разделения ответственности, что отражается в выделении четырёх горизонтальных уровней: интерфейса пользователя, программного ядра, сцены и рендеринга, а также физического моделирования. Рисунок~\ref{fig:modeuler_architecture} иллюстрирует взаимодействие этих уровней.

\begin{figure}[h]
\centering
\begin{tikzpicture}[
  font=\small,
  layer/.style={
    draw, rounded corners=4pt,
    minimum width=11cm, minimum height=1.1cm,
    align=center, text width=10.5cm
  },
  grouplabel/.style={font=\footnotesize\itshape, text=gray},
  arr/.style={-{Stealth[length=6pt]}, thick},
]

% Layers (bottom to top)
\node[layer, fill=blue!10]   (gpu)     at (0,0)   {GPU-ускоренный физический движок\\{\footnotesize Vulkan Compute, SSBO, вычислительные шейдеры}};
\node[layer, fill=green!10]  (scene)   at (0,1.5) {Слой сцены и рендеринга\\{\footnotesize VulkanSceneGraph v1.1.14, vsgXchange, urdfdom}};
\node[layer, fill=yellow!10] (core)    at (0,3.0) {Программное ядро (C~ABI)\\{\footnotesize modeuler.h, управление сценой, Vulkan-устройство}};
\node[layer, fill=orange!10] (ui_win)  at (-2.75,4.5) {WPF / .NET~10\\{\footnotesize Windows}};
\node[layer, fill=orange!10, minimum width=5cm, text width=4.5cm] (ui_gtk)  at (2.75,4.5) {GTK~4 / C\\{\footnotesize Linux}};

% Arrows
\draw[arr] (gpu.north)   -- (scene.south);
\draw[arr] (scene.north) -- (core.south);
\draw[arr] (core.north) -- ++(0,0.55) -| (ui_win.south);
\draw[arr] (core.north) -- ++(0,0.55) -| (ui_gtk.south);

% Side labels
\node[grouplabel, right=0.3cm of gpu]   {§\,3.3};
\node[grouplabel, right=0.3cm of scene] {VSG};
\node[grouplabel, right=0.3cm of core]  {§\,3.1.3};

\end{tikzpicture}
\caption{Многоуровневая архитектура программного комплекса «Модейлер»}
\label{fig:modeuler_architecture}
\end{figure}

\textbf{Уровень пользовательского интерфейса.} Для обеспечения кроссплатформенности реализованы два независимых приложения верхнего уровня. Приложение для Windows построено на основе .NET~10 с использованием Windows Presentation Foundation (WPF) и библиотек Fluent.Ribbon и AvalonDock, обеспечивающих профессиональный пользовательский интерфейс с вкладками, плавающими панелями и лентой инструментов. Приложение для Linux реализовано на C с использованием GTK~4. Оба варианта взаимодействуют с ядром исключительно через C~ABI, что позволяет независимо развивать каждый из интерфейсов.

\textbf{Уровень программного ядра} предоставляет публичный стабильный интерфейс в виде заголовочного файла~\texttt{modeuler.h}, полностью совместимого со стандартом языка~C. Такое решение обеспечивает двоичную совместимость (ABI stability): интерфейсные типы используют только примитивные типы данных и непрозрачные указатели (\texttt{modeuler\_scene\_t}, \texttt{modeuler\_window\_t}), а все детали реализации на C++ скрыты за непрозрачным указателем. Внутри ядра располагается менеджер Vulkan-устройства — единственный экземпляр~\texttt{vsg::Instance} и связанное с ним~\texttt{vsg::Device}, создаваемые при первом обращении.

\textbf{Уровень сцены и рендеринга} реализован на основе VulkanSceneGraph~(VSG)~--- высокопроизводительной библиотеки визуализации, построенной непосредственно поверх Vulkan. VSG организует трёхмерную сцену в виде ориентированного ациклического графа (DAG), узлы которого соответствуют геометрическим преобразованиям, меш-данным и материалам. Для загрузки форматов 3D-геометрии используется плагин~vsgXchange; для парсинга описаний роботов --- библиотека~urdfdom, разработанная в рамках экосистемы ROS.

\textbf{Физический движок} --- новый компонент, разрабатываемый в рамках данной диссертационной работы. Он занимает нижний уровень архитектуры и взаимодействует со слоем сцены в обоих направлениях: получает от него структуру кинематической цепочки модели и передаёт обновлённые трансформации тел на каждом шаге симуляции.

\subsection{Кинематическая модель робота}

Описание робота в формате URDF задаёт дерево звеньев (\textit{links}) и сочленений (\textit{joints}). Каждое звено несёт информацию о визуальной геометрии, коллизионной геометрии и инерционных характеристиках; каждое сочленение определяет тип кинематической пары, её ось и пределы перемещения. Система поддерживает следующие типы сочленений:

\begin{itemize}
  \item \texttt{REVOLUTE} — вращательная пара с ограниченным углом;
  \item \texttt{CONTINUOUS} — вращательная пара без ограничения угла;
  \item \texttt{PRISMATIC} — поступательная пара;
  \item \texttt{FLOATING} — 6-DOF сочленение (свободное тело);
  \item \texttt{PLANAR} — движение в плоскости;
  \item \texttt{FIXED} — жёсткое соединение без степеней свободы.
\end{itemize}

После загрузки URDF-файла физический движок строит внутреннее представление кинематической цепочки: каждое звено становится объектом типа \texttt{RigidBody} (или статическим препятствием), а каждое сочленение --- ограничением (\textit{constraint}) в системе уравнений динамики.

\subsection{Программный интерфейс ядра}
\label{sec:core_api}

Программный интерфейс ядра спроектирован по принципу минимальной поверхности: пользовательский код работает только с непрозрачными дескрипторами и узким набором функций. Это позволяет изменять внутреннюю реализацию без нарушения обратной совместимости.

\begin{lstlisting}[language=C, caption={Ключевые функции публичного C-интерфейса «Модейлера»}]
/* Инициализация и завершение работы */
modeuler_result_t modeuler_init(void);
void              modeuler_deinit(void);

/* Управление сценой */
modeuler_scene_t  modeuler_scene_create(void);
void              modeuler_scene_destroy(modeuler_scene_t scene);

/* Загрузка модели и создание физических тел */
modeuler_model_t  modeuler_model_import(modeuler_scene_t scene,
                                        const char *urdf_path);

/* Управление окном отображения */
modeuler_window_t modeuler_window_attach_native(
                      modeuler_scene_t scene,
                      void *native_handle,
                      modeuler_native_type_t type);

/* Шаг симуляции */
modeuler_result_t modeuler_physics_step(modeuler_scene_t scene,
                                        double dt);

/* Рендеринг кадра */
modeuler_result_t modeuler_run_frame(modeuler_window_t window);
\end{lstlisting}

Функция \texttt{modeuler\_physics\_step} вызывается управляющим циклом приложения с фиксированным шагом~$\Delta t$ (типично 1--2\,мс для реального времени). Результаты шага симуляции (обновлённые трансформации тел) немедленно применяются к VSG-узлам, после чего вызов~\texttt{modeuler\_run\_frame} генерирует новый кадр изображения.

Привязки для управляемых языков формируются инструментом~ClangSharp, который разбирает заголовочный файл~\texttt{modeuler.h} и генерирует код P/Invoke на языке~C\# для интеграции с WPF-приложением.

\subsection{Взаимодействие с ROS~2}

Интеграция с ROS~2 реализуется на двух уровнях. На уровне форматов данных: «Модейлер» использует URDF~--- официальный формат описания роботов в экосистеме ROS, а система координат при загрузке трансформаций соответствует конвенции фреймов~TF2. На уровне обмена данными: физический движок публикует состояние сочленений (\textit{joint states}), одометрию и показания сенсоров через топики ROS~2, что позволяет использовать «Модейлер» в качестве имитатора для управляющего программного обеспечения, работающего в штатной среде ROS~2 без каких-либо изменений кода.
